\section{结论}
\label{sec:conclusion-zh}

本文提出 DVCL，一种面向鲁棒异构图神经网络的双视图对比学习框架。DVCL 将净化语义拓扑视图与特征诱导视图结合，并通过对称跨视图对比目标对齐两个视图的节点表示。该设计在保留异构语义信息的同时，降低了模型对可能已被污染的图拓扑的单一依赖。

本文的实验框架只使用准确率指标，并要求所有方法采用完全相同的数据划分、攻击产物、目标节点、随机种子和调参规则。在当前统一 ACM 协议下，预备证据表明，DVCL 在 PRBCD、HetePRBCD 和 FGSM 风格结构目标攻击下的准确率高于本文的 HSeCo 复现。该结论有意保持克制：只有在完成 HSeCo checkpoint 审计、多个完整数据集以及 RoHe、HeteGuard 等外部基线后，论文才能给出更广泛的先进性结论。

当前评估主要关注结构扰动。由于 DVCL 同时依赖特征诱导视图，还需要通过特征攻击以及联合攻击拓扑和属性的自适应攻击来刻画完整鲁棒性边界。未来工作还可以将 DVCL 扩展到动态异构图和其他攻击面。
